% =====================================================================
% SHARD-ID: LB-02-MPS-SETTING
% SHARD-TITLE: The setting: uniform matrix-product states, windows, and
%              excitations
% SHARD-SOURCES: definitions.md D1, D5, D9, D11, D12 (status header: written
%              2026-08-25, FROZEN at the 2026-08-26 freeze after the Corner A
%              L6 loop converged, theory/verdicts/corner-a-r3.md PASS);
%              notation.md; HANDOFF_MPS_SOFT_THEOREM.md 1.1, 1.3;
%              theory/corner-a.md, theory/corner-a-goldstone.md,
%              theory/corner-a-kinks.md; theory/checks/corner_a_check.py
%              (checks C4, C5, C10, C11)
% SHARD-CLAIMS: D1, D5, D9, D11, D12
% =====================================================================

% --- shard-local semantic macros (mirror notation.md; \providecommand so
% --- that a sibling shard may declare the same symbol without a clash) ---
\providecommand{\Aloc}{\mathfrak{A}_{\mathrm{loc}}}      % local subalgebra
\providecommand{\Utens}{\mathcal{U}}                     % symmetry on a tensor
\providecommand{\Vac}{\Omega_{\mathrm{vac}}}             % vacuum label set
\providecommand{\Nnull}{\mathcal{N}}                     % null-direction map
\providecommand{\Ewin}{\mathfrak{E}}                     % window CP map
\providecommand{\lamt}{\tilde{\lambda}}                  % honest decay rate
\providecommand{\Cbd}{C_{\partial}}                      % window-norm constant
\providecommand{\gaugerem}{\mathfrak{B}}                 % gauge remainder
\providecommand{\Ept}{\mathcal{E}}                       % endpoint space
\providecommand{\dcoset}{\mathfrak{d}}                   % double-coset label
\providecommand{\Fc}{\mathfrak{F}_{c}}                   % finitely supported
\providecommand{\Fec}{\mathfrak{F}_{\mathrm{ec}}}        % eventually constant
\providecommand{\Fell}{\mathfrak{F}_{\ell^{1}}}          % finite variation
\providecommand{\Aeff}{\A_{\mathrm{eff}}}                % effective asg group
\providecommand{\hsub}{\mathfrak{h}}                     % unbroken subalgebra
\providecommand{\msub}{\mathfrak{m}}                     % broken complement
\providecommand{\liealg}{\mathfrak{g}}                   % Lie algebra of G

\section{The setting: uniform matrix-product states, windows, and excitations}
\label{sec:mps-setting}

A uniform matrix-product state is a translation-invariant state of an infinite
spin chain built from a single small tensor. One fixes a physical dimension
$d$ and a \emph{bond dimension} $\chi$, and one array $A = (A^{s})_{s=1}^{d}$
of $\chi \times \chi$ matrices. Reading a physical configuration
$s_{1}s_{2}\cdots s_{w}$ of $w$ consecutive sites as the ordered matrix word
$A^{s_{1}}A^{s_{2}}\cdots A^{s_{w}}$, and closing that word against two fixed
matrices, produces the amplitudes of a state on any window; the state on the
whole chain is the compatible family of all such windows. All correlations of
that state are governed by a single completely positive map on
$\chi \times \chi$ matrices, the \emph{transfer matrix}
$E(Y) = \sum_{s} A^{s} Y (A^{s})^{\dagger}$. When $E$ has a unique, simple
eigenvalue of largest modulus --- the \emph{injective} case --- the resulting
state is pure, translation invariant, and exponentially clustering, with
correlation length set by the gap between the leading eigenvalue $1$ and the
rest of the spectrum. Injectivity is the whole reason such states are
tractable: it makes the state unique, it makes the symmetry data of
\cref{sec:symmetry-noether} unique, and it makes long matrix words span the
full matrix algebra, which is what lets a bond insertion be recovered from the
window vector it produces.

The campaign is about what a symmetry does when it is applied to a
\emph{finite region} rather than the whole chain, so nothing may be said only
about the infinite chain. This is why the basic object below is not the state
but the \emph{window vector}: a finite interval $\Lambda = [a,b]$, the ordered
product of the tensors sitting on it, arbitrary invertible matrices inserted on
any of its bonds, and two boundary vectors closing the word. Everything the
campaign proves --- the truncated-symmetry identity, the bond implementers, the
soft limits --- is first an exact identity between window vectors, and only
afterwards a statement about states. Two technical points earn their keep
immediately. First, the $\abs{\Lambda}+1$ bonds available for insertion must
include the two \emph{edge} bonds $(a-1|a)$ and $(b|b+1)$; without them the
truncated-symmetry identity is simply false on the window $\Lambda = R$ that it
is about. Second, the decay of $E^{m}$ towards its fixed point must be quoted
honestly: $E$ may carry Jordan blocks at the subleading modulus, so the correct
statement carries a rate $\lamt$ strictly above the transfer gap $\lambda_{E}$,
never $\lambda_{E}$ itself. Both points are recorded in the definitions below
because both were live errors at some stage of the campaign.

\subsection{Vacua: uniform injective matrix-product states}

The first definition fixes the arena --- the algebra of observables on the
chain --- and the class of states the whole campaign treats as ``vacua''. The
content of the injectivity clause is that the transfer matrix has a spectral
gap, so that the state clusters exponentially and long matrix words are
unconstrained.

\begin{definition}[Uniform injective matrix-product vacuum (D1(a)--(d))]
\label{def:injective-mps}
\emph{(a) The chain.} Fix $d \in \N$ and let the lattice be $\Z$. The
\emph{quasi-local algebra} is
\begin{equation}
  \alg := \bigotimes_{x \in \Z} M_{d}(\C),
\end{equation}
the C*-inductive limit of the window algebras
$\alg_{\Lambda} := \bigotimes_{x \in \Lambda} M_{d}(\C)$ over finite
$\Lambda \subset \Z$; the norm-dense \emph{local subalgebra} is
$\Aloc := \bigcup_{\Lambda} \alg_{\Lambda}$. Translation by $x$ sites is the
automorphism $\tau_{x} \in \mathrm{Aut}(\alg)$. A \emph{state} is a normalised
positive linear functional on $\alg$; states are compared in the weak-$*$
topology, i.e.\ pointwise on $\Aloc$. For a state $\varrho$ we write
$(\Hilb_{\varrho}, \pi_{\varrho}, \Omega_{\varrho})$ for its GNS triple.

\emph{(b) Tensor and transfer matrix.} A \emph{uniform matrix-product tensor}
of bond dimension $\chi$ is $A = (A^{s})_{s=1}^{d}$ with
$A^{s} \in M_{\chi}(\C)$. Its \emph{transfer matrix} is the completely
positive map $E : M_{\chi} \to M_{\chi}$,
\begin{equation}
  E(Y) \; := \; \sum_{s} A^{s} \, Y \, (A^{s})^{\dagger},
\end{equation}
equivalently the $\chi^{2} \times \chi^{2}$ matrix
$\sum_{s} A^{s} \otimes \overline{A^{s}}$.

\emph{(c) Injectivity, canonical form, gap.} $A$ is \emph{injective}
(equivalently \emph{normal}) if $E$ has spectral radius $1$, the eigenvalue $1$
is the unique eigenvalue of modulus $1$, it is non-degenerate, and the left and
right eigenmatrices $l, r$ --- i.e.\ $E^{*}(l) = l$ and $E(r) = r$ --- are
positive definite. We fix the \emph{canonical form}
\begin{equation}
  l = \mathbb{1}, \qquad r > 0, \qquad \operatorname{tr} r = 1 ,
\end{equation}
after which the residual matrix-product gauge freedom is
$A^{s} \mapsto Y^{-1} A^{s} Y$ with $Y$ unitary commuting with $r$, together
with an overall phase. The \emph{transfer gap} is
\begin{equation}
  \lambda_{E} \; := \; \max \bigl\{\, \abs{\mu} \;:\; \mu \in \spec(E), \;
  \mu \neq 1 \,\bigr\} \; < \; 1 ,
  \qquad \xi_{c} := -1/\log \lambda_{E} ,
\end{equation}
with $\xi_{c}$ the correlation length. Because $E$ may have Jordan blocks at
modulus $\lambda_{E}$, the honest decay statement is: \emph{for every
$\lamt \in (\lambda_{E},1)$ there is $C_{\lamt} < \infty$ with}
\begin{equation}
  \norm{E^{m} - P} \;\le\; C_{\lamt} \, \lamt^{\,m}
  \qquad \text{for all } m \ge 0,
  \label{eq:honest-decay}
\end{equation}
where $P(Y) = \operatorname{tr}(Y)\, r$ is the spectral projection onto the
fixed point. A bare $O(\lambda_{E}^{m})$ is \emph{false} in general: it needs
the polynomial prefactor $m^{p}$ with $p$ equal to the Jordan block size minus
one. Every rate quoted anywhere in this labbook is in the $\lamt$ form.
Injectivity is equivalent to the existence of an $n_{0}$ with
\begin{equation}
  \operatorname{span}\bigl\{\, A^{s_{1}} \cdots A^{s_{n}} \;:\;
  s_{1} \ldots s_{n} \,\bigr\} \;=\; M_{\chi}(\C)
  \qquad \text{for all } n \ge n_{0} .
  \label{eq:word-spanning}
\end{equation}

\emph{(d) The state.} For injective $A$ in canonical form, $\omega_{A}$ is the
unique state on $\alg$ with, for $O \in \alg_{[1,w]}$,
\begin{equation}
  \omega_{A}(O) \;=\; \sum_{s,s'} \bra{s'} O \ket{s} \,
  \operatorname{tr}\Bigl[\, (A^{s'_{w}})^{\dagger} \cdots (A^{s'_{1}})^{\dagger}
  \cdot \mathbb{1} \cdot A^{s_{1}} \cdots A^{s_{w}} \cdot r \,\Bigr],
\end{equation}
extended by translation invariance. The state $\omega_{A}$ is pure,
translation invariant, and exponentially clustering with rate $\lambda_{E}$
(in the sense of \eqref{eq:honest-decay}). We write
$(\Hilb_{A}, \pi_{A}, \Omega_{A})$ for its GNS triple and suppress $\pi_{A}$
when harmless.
\end{definition}

\provenance{D1(a)--(d)}%
{\texttt{definitions.md} D1, clauses (a)--(d); FROZEN at the 2026-08-26 freeze}%
{---}%
{\texttt{theory/verdicts/corner-a-r3.md} (PASS)}

\subsection{Windows, decorations, and window vectors}

A decoration is a window with things put on it: possibly different tensors at
some sites, possibly invertible matrices squeezed onto some bonds, and two
boundary vectors at the ends. This is the working object of the campaign
because every symmetry statement below turns out to be an identity between a
plain window vector and a decorated one. The edge-bond clause is not
cosmetic: it is what makes ``the symmetry applied to $R$ acts only on the two
boundary bonds of $R$'' a statement about the window $R$ itself rather than
about some larger window.

\begin{definition}[Window vectors, decorations, and decorated states (D1(e))]
\label{def:window-vector}
For a finite window $\Lambda = [a,b]$, a \emph{decoration} is a tuple
$T = (T_{a},\dots,T_{b})$ of tensors $T_{x} = (T_{x}^{s})$, together with
\emph{bond insertions} $M_{x|x+1} \in GL(\chi)$ on any subset of the
$\abs{\Lambda}+1$ bonds
\begin{equation}
  (a-1|a), \quad (a|a+1), \quad \dots, \quad (b|b+1)
\end{equation}
--- \textbf{including the two edge bonds} --- and boundary vectors
$b_{l}, b_{r} \in \C^{\chi}$. The \emph{window vector} is
\begin{equation}
  \ket{\psi_{\Lambda}(T; b_{l}, b_{r})} \; := \!\!
  \sum_{s_{a} \ldots s_{b}} \!\!
  \bra{b_{l}} M_{a-1|a} \, T_{a}^{s_{a}} \, M_{a|a+1} \cdots
  T_{b}^{s_{b}} \, M_{b|b+1} \ket{b_{r}} \; \ket{s_{a} \ldots s_{b}},
  \label{eq:window-vector}
\end{equation}
absent factors being $\mathbb{1}$. An edge insertion is equivalently a
redefinition of the boundary vector,
\begin{equation}
  b_{l} \mapsto M_{a-1|a}^{\dagger} \, b_{l},
  \qquad
  b_{r} \mapsto M_{b|b+1} \, b_{r} ,
\end{equation}
and is included so that identities whose insertions land on $\partial \Lambda$
are statements about the same window rather than about a larger one.

There is a constant $\Cbd = \Cbd(A, b_{l}, b_{r}) < \infty$ with
\begin{equation}
  \bigl\Vert \ket{\psi_{\Lambda}(T; b_{l}, b_{r})} \bigr\Vert
  \;\le\; \Cbd \prod_{M} \norm{M}
  \label{eq:cbd}
\end{equation}
for every $\Lambda$ and every decoration differing from the uniform $A$ on
boundedly many sites --- uniformly in $\abs{\Lambda}$, because the contraction
is of the form $(b_{l} \otimes \bar{b}_{l} \vert E^{\abs{\Lambda}} \vert
b_{r} \otimes \bar{b}_{r})$ and $E^{m}$ is power-bounded by
\cref{def:injective-mps}(c).

A \emph{decorated state} $\omega_{A}[T]$ is the state on $\alg$ obtained from a
decoration that differs from the uniform $A$ only on finitely many sites and
bonds, by the prescription of \cref{def:injective-mps}(d) with the modified
string and with $l, r$ as the two infinite environments; it is well defined by
the $\lamt$-estimate \eqref{eq:honest-decay}. We write $\omega_{A}^{M@b}$ for
the state decorated by a single bond insertion $M$ on the bond $b$.
\end{definition}

\begin{remark}[Why the edge bonds are in the definition]
\label{rem:edge-bond-repair}
The clause admitting insertions on $(a-1|a)$ and $(b|b+1)$ was added in
revision r2 of the Corner~A loop. The earlier window vector carried interior
bonds only, which made the earlier statement of the truncated-symmetry
identity false for the window $\Lambda = R$ --- the two surviving virtual
operators then lie outside the expression altogether. The measured
discrepancy for the $\chi = 2$ Pauli tensor was $1.658$; admitting edge
insertions restores exactness to $0.0$.
\end{remark}

\provenance{D1(e)}%
{\texttt{definitions.md} D1(e); repair recorded at \texttt{theory/corner-a.md}
 step (1.3), sub-step (2.7)}%
{\texttt{theory/checks/corner\_a\_check.py}, checks C2, C2b}%
{\texttt{theory/verdicts/corner-a-r2.md} objection 1;
 \texttt{theory/verdicts/corner-a-r3.md}}

A kink state has a different vacuum on each half-line, so it is not a
decoration in the above sense at all: it differs from any uniform tensor on a
whole half-line. It is defined instead by its finite-window restrictions, and
those restrictions are matrix-environment contractions rather than window
vectors --- a distinction that cost the campaign one round of corrections.

\begin{definition}[Two-sided (half-line) decorations (D1(e$'$))]
\label{def:two-sided-decoration}
A \emph{two-sided decoration} carries an injective canonical-form tensor
$A_{\alpha}$ on $(-\infty, m]$ and an injective canonical-form tensor
$A_{\beta}$ on $[m+1, \infty)$ (possibly $\alpha = \beta$), plus finitely many
further site and bond modifications. For a finite window $W \ni m$ let
$P_{W}(s)$ be the ordered product of the window tensors, with the
modifications inserted, for the physical index string $s$, and define the
\emph{unnormalised completely positive contraction}
\begin{align}
  \tilde{\varrho}_{W}
  &:= \sum_{s,s'} \operatorname{tr}\bigl[\, l_{\alpha} \, P_{W}(s) \,
     r_{\beta} \, P_{W}(s')^{\dagger} \,\bigr] \; \ket{s}\bra{s'},
  \label{eq:two-sided-contraction} \\[2pt]
  \omega_{\alpha|\beta}^{(m)}[T](O)
  &:= \frac{\operatorname{tr}[\tilde{\varrho}_{W} O]}
          {\operatorname{tr}[\tilde{\varrho}_{W}]}
  \qquad \text{for } O \in \alg_{W}.
\end{align}
\end{definition}

\begin{lemma}[The two-sided prescription defines a state]
\label{lem:two-sided-welldefined}
$\tilde{\varrho}_{W} \ge 0$, $\operatorname{tr} \tilde{\varrho}_{W} > 0$, and
the functionals are consistent under $W \subseteq W'$; hence they define a
unique state on $\alg$.
\end{lemma}

\begin{proof}[Idea of proof]
\emph{Positivity.} For $c \in \C^{d^{\abs{W}}}$ put
$T := \sum_{s} c_{s} P_{W}(s)$. Then
\begin{equation}
  \bra{c} \tilde{\varrho}_{W} \ket{c}
  \;=\; \operatorname{tr}\bigl[\, l_{\alpha} T r_{\beta} T^{\dagger} \,\bigr]
  \;=\; \operatorname{tr}\Bigl[\,
        \bigl(l_{\alpha}^{1/2} T r_{\beta}^{1/2}\bigr)
        \bigl(l_{\alpha}^{1/2} T r_{\beta}^{1/2}\bigr)^{\dagger} \,\Bigr]
  \;\ge\; 0 ,
\end{equation}
using $l_{\alpha}, r_{\beta} > 0$ from \cref{def:injective-mps}(c).
\emph{Nondegeneracy.}
$\operatorname{tr} \tilde{\varrho}_{W}
 = \operatorname{tr}[\, l_{\alpha} \, \Ewin_{W}(r_{\beta}) \,] > 0$
with $\Ewin_{W}$ the window completely positive map, since
$l_{\alpha}, r_{\beta} > 0$ and $\Ewin_{W}$ is faithful on positive definite
inputs. \emph{Consistency.} Enlarging $W$ by one site on the left tail
applies $\Ewin_{A_{\alpha}}^{*}$ to $l_{\alpha}$, which fixes it; on the right
tail it applies $\Ewin_{A_{\beta}}$ to $r_{\beta}$, which fixes it.
\end{proof}

\begin{remark}[A corrected argument]
The r2 text justified positivity by saying that each entry is
$\bra{v} \cdot \ket{v}$ for a finite-window vector $v$. That is wrong: with
full-rank matrix environments the restriction is a positive completely
positive contraction, and becomes a vector expression only after purifying the
environments --- not a rank-one boundary-vector window vector. The argument
above is the corrected one. Independently, the weak-$*$ limit constructed in
the Corner~A endpoint theorem supplies positivity, so
\cref{lem:two-sided-welldefined} is a convenience rather than a load-bearing
hypothesis.
\end{remark}

\provenance{D1(e$'$)}%
{\texttt{definitions.md} D1(e$'$), with the well-definedness lemma stated and
 proved there}%
{\texttt{theory/checks/corner\_a\_check.py}, check C11 (Hermiticity $0$,
 minimal eigenvalue $-1.4 \cdot 10^{-16}$, positive trace)}%
{\texttt{theory/verdicts/corner-a-r2.md} objection 8 (conceded);
 \texttt{theory/verdicts/corner-a-r3.md}}

\subsection{Excitations: the ansatz and its gauge freedom}

One-particle states are built by replacing the vacuum tensor at one site by a
different tensor $B$ and superposing over the position of that replacement with
a momentum phase. If the tensors to the left and to the right of $B$ are the
same, the result is a local excitation over a single vacuum; if they are two
different vacua, it is a kink. The parametrisation is redundant, and the
redundancy --- the \emph{null directions} --- is the technical heart of the
soft theorem: at zero momentum the Goldstone tensor of
\cref{def:goldstone-tensor} becomes a null direction, which is exactly the
statement that the zero-momentum Goldstone mode is a global symmetry rotation
and therefore carries no independent excitation.

\begin{definition}[Excitation ansatz and kink sectors (D5(a))]
\label{def:excitation-ansatz}
For injective canonical-form tensors $A_{\alpha}, A_{\beta}$, a tensor
$B = (B^{s})$ with $B^{s} \in M_{\chi}$, and $k \in (-\pi,\pi]$, the
\emph{excitation-ansatz vector} is defined on windows by
\begin{equation}
  \ket{\Phi_{k}(B; A_{\alpha}, A_{\beta})}_{\Lambda}
  \; := \; \sum_{n \in \Lambda} e^{ikn} \,
  \ket{\psi_{\Lambda}\bigl( \ldots A_{\alpha} A_{\alpha} \, B \,
       A_{\beta} A_{\beta} \ldots ; \, b_{l}, b_{r} \bigr)},
  \label{eq:excitation-ansatz}
\end{equation}
with $B$ at site $n$, $A_{\alpha}$ at all sites $< n$ and $A_{\beta}$ at all
sites $> n$. For $\alpha = \beta$ we write $\ket{\Phi_{k}(B)}$ and call it the
\emph{trivial} (local) sector; for $\alpha \neq \beta$ it is the \emph{kink}
(topological) sector. On the infinite chain these are $\delta$-normalised
generalised vectors,
\begin{equation}
  \braket{\Phi_{k}(B) \vert \Phi_{k'}(B')}
  \;=\; 2\pi \, \delta(k-k') \cdot B^{\dagger} N_{k} B' ,
  \label{eq:delta-normalisation}
\end{equation}
with $N_{k}$ the Gram form of the ansatz at momentum $k$.
\end{definition}

\provenance{D5(a)}%
{\texttt{definitions.md} D5(a); convention matched against the local \TeX{}
 ground truth \texttt{refs/arxiv-1103.2286} (\texttt{dispersionrelation\_final.tex})
 and \texttt{refs/arxiv-1810.07006} (\texttt{p5\_excitations.tex}, Eq.~(kink))}%
{---}%
{\texttt{theory/verdicts/corner-a-r3.md}}

\begin{definition}[Null directions of the ansatz (D5(b))]
\label{def:null-directions}
For $X \in M_{\chi}$ set
\begin{equation}
  \Nnull_{k}^{\alpha\beta}(X) \; := \; e^{ik} A_{\alpha} X - X A_{\beta},
  \qquad\text{i.e.}\qquad
  \Nnull_{k}(X)^{s} = e^{ik} A_{\alpha}^{s} X - X A_{\beta}^{s} .
  \label{eq:null-directions}
\end{equation}
This is the gauge freedom $B \simeq B + e^{ik} A X - X A$ of the ansatz.
The rank of $X \mapsto \Nnull_{k}(X)$ is $\chi^{2}$ for $k \neq 0$ and
$\chi^{2} - 1$ for $k = 0$ with $\alpha = \beta$: the rank \emph{drops} by one
as $k \to 0$.
\end{definition}

\begin{remark}[The vanishing of a null direction is a limit statement, not a
finite-window identity]
\label{rem:null-not-exact}
The identity $\ket{\Phi_{k}(\Nnull_{k}(X))} = 0$ is \textbf{false as a
finite-window identity}. On $\Lambda = [a,b]$ the reindexing that would prove
it leaves exactly two boundary terms,
\begin{equation}
  e^{ik(b+1)} \ket{\psi_{\Lambda}; \, X@(b|b+1)}
  \; - \; e^{ika} \ket{\psi_{\Lambda}; \, X@(a-1|a)} ,
  \label{eq:null-boundary-terms}
\end{equation}
which do not vanish for generic $b_{l}, b_{r}$: the measured norm is $1.77$ for
the $\chi = 2$ Pauli tensor at $k = 0.37$ (an earlier measurement with
different boundary vectors gave $0.591$). The vanishing statement holds only
in the limits specified in \cref{def:gauge-limits}. Shards of the first
revision asserted the exact identity and were wrong; the exact identity, with
its boundary terms, is the summation-by-parts lemma of the Goldstone shard.
\end{remark}

\provenance{D5(b)}%
{\texttt{definitions.md} D5(b); exact identity with boundary terms at
 \texttt{theory/corner-a-goldstone.md} step (1.5) (Lemma SBP); rank statement
 ibid.; ground truth \texttt{refs/arxiv-1103.2286} and
 \texttt{refs/arxiv-1810.07006} Eq.~(gaugeexcitations)}%
{\texttt{theory/checks/corner\_a\_check.py}}%
{\texttt{theory/verdicts/corner-a-r2.md} (warning conceded);
 \texttt{theory/verdicts/corner-a-r3.md}}

\subsection{Superselection: vacuum and kink sectors}

The next definition says what it means for a state to ``look like vacuum
$\alpha$ far to the left and vacuum $\beta$ far to the right''. The
factorised form of the boundary condition --- with a spectator observable $D$
inserted --- is what makes the label stable under normal perturbations, and
hence a genuine superselection label rather than a property of one particular
state. It refers forward to the covariant vacuum family of
\cref{def:covariant-vacuum}.

\begin{definition}[Vacuum and kink superselection sectors (D9(a),(b))]
\label{def:kink-sector}
Let $\{\omega_{\alpha}\}_{\alpha \in \Vac}$ be a $G$-covariant vacuum family
(\cref{def:covariant-vacuum}). For $\alpha, \beta \in \Vac$ the \emph{kink
sector} $\Kk_{\alpha\beta}$ is the set of states $\varrho$ on $\alg$ obeying
the \emph{factorised} boundary conditions at infinity,
\begin{equation}
  \lim_{n \to -\infty} \varrho\bigl(D \cdot \tau_{n}(O)\bigr)
    = \varrho(D) \, \omega_{\alpha}(O),
  \qquad
  \lim_{n \to +\infty} \varrho\bigl(D \cdot \tau_{n}(O)\bigr)
    = \varrho(D) \, \omega_{\beta}(O),
  \label{eq:kink-boundary-conditions}
\end{equation}
for all $D, O \in \Aloc$. Taking $D = \mathbb{1}$ gives the plain conditions
$\varrho(\tau_{n}(O)) \to \omega_{\alpha/\beta}(O)$. The sector
$\Kk_{\alpha\alpha}$ is the \emph{vacuum sector of $\alpha$}; it contains
$\omega_{\alpha}$, every state $\omega_{\alpha} \circ \Ad(W^{\dagger})$ with
$W \in \alg$ unitary, and every decorated state $\omega_{A_{\alpha}}[T]$ of
\cref{def:window-vector}. A state in $\Kk_{\alpha\beta}$ with $\alpha \neq \beta$
is a \emph{kink state} (equivalently a domain-wall or topological state).

\emph{Disjointness.} States in $\Kk_{\alpha\beta}$ and
$\Kk_{\alpha'\beta'}$ with $(\alpha,\beta) \neq (\alpha',\beta')$ are disjoint,
i.e.\ mutually non-normal. Hence ``superselection sector'' is used here in the
precise sense: the label $(\alpha,\beta)$ is a classical observable at
infinity, unchangeable by any element of $\alg$.
\end{definition}

\provenance{D9(a),(b)}%
{\texttt{definitions.md} D9(a),(b); stability of the factorised label and
 disjointness proved at \texttt{theory/corner-a-kinks.md} step (1.8), with
 rate $\lambda_{E}$ for every decorated matrix-product state at sub-step
 (1.8)(iii)}%
{---}%
{\texttt{theory/verdicts/corner-a-r3.md}}

The next object is the set of states obtained from one vacuum by inserting a
single invertible matrix on one bond. It is defined here as a bare set with a
bare action; \emph{that} the action descends, and what its orbits are, are
theorems of the Corner~A shard, restated below only for the reader's
convenience.

\begin{definition}[Endpoint space and endpoint action (D9(c), definition only)]
\label{def:endpoint-space}
Fix a bond $b$ and a vacuum label $\alpha$. The \emph{endpoint space} is
\begin{equation}
  \Ept_{b}^{\alpha}
  \; := \; \bigl\{\, \omega_{A_{\alpha}}^{M@b} \;:\; M \in GL(\chi) \,\bigr\},
\end{equation}
and the \emph{endpoint action} of $G$ on representatives is
\begin{equation}
  g \star M \; := \; V_{\alpha}(g) \, M .
\end{equation}
Nothing beyond this is definitional.
\end{definition}

\begin{remark}[What the endpoint theorem adds]
\label{rem:endpoint-torsor}
The following are proved in the Corner~A endpoint theorem and are recorded
here only because they are what makes \cref{def:endpoint-space} interesting.
(i) $\omega^{M@b} = \omega^{M'@b}$ if and only if $M' \in \C^{\times} M$;
hence $[M] \mapsto \omega^{M@b}$ is a canonical bijection
$PGL(\chi) \to \Ept_{b}^{\alpha}$, and $\Ept_{b}^{\alpha}$ is a
$PGL(\chi)$-torsor. (ii) The action $\star$ descends to $\Ept_{b}^{\alpha}$
and acts there as left translation by
$\rho_{\alpha}(g) = [V_{\alpha}(g)] \in PGL(\chi)$, a genuine $G$-action with
kernel $N_{\alpha}$; the orbit of $\omega_{\alpha}$ is
$\rho_{\alpha}(G) \cong \Aeff = G/N_{\alpha}$, on which $\Aeff$ acts simply
transitively. (iii) The lift of $\rho_{\alpha}$ to the padded window space
furnished by \cref{def:bond-implementer} is a \emph{projective} action with
multiplier $\omega_{\alpha}$; it exists always, including when
$[\omega_{\alpha}] \neq 0$. What the class $[\omega_{\alpha}]$ obstructs is
\emph{removing} the multiplier, i.e.\ lifting $\rho_{\alpha}$ to an honest
homomorphism $G \to U(\chi)$. The earlier phrasing, that
$[\omega_{\alpha}]$ obstructs ``lifting the state action to the window
action'', is wrong and has been retracted.

Two sanity checks. For $\chi = 1$, or whenever $V_{\alpha}$ is a scalar for
every $g$, one has $N_{\alpha} = G$ and a one-point orbit. For the AKLT chain
with $G = \Z_{2} \times \Z_{2}$ and $V_{\alpha} \in \{\mathbb{1}, X, Z, XZ\}$
one has $N_{\alpha} = \{e\}$ and a four-point orbit.

When it is said that the asymptotic symmetry ``relabels charge superselection
sectors'', that is asserted about this torsor and about the module structure of
$\C^{\chi}$ over the charge algebra of \cref{def:charge-algebra} --- never
about $\Kk_{\alpha\beta}$, which is trivial in the unbroken case.
\end{remark}

\provenance{D9(c),(c$'$)}%
{\texttt{definitions.md} D9(c) (definition) and D9(c$'$) (corollaries);
 proved at \texttt{theory/corner-a.md} step (1.4), sub-steps (d2) and (f)}%
{\texttt{theory/checks/corner\_a\_check.py}}%
{\texttt{theory/verdicts/corner-a-r2.md} objections 2 and 7;
 \texttt{theory/verdicts/corner-a-r3.md}}

Finally, when the symmetry is broken there is a question of which pairs of
vacua are genuinely different once the global symmetry is quotiented out. The
naive answer --- the difference $g_{L} g_{R}^{-1}$ --- is wrong for nonabelian
groups; the correct invariant is a double coset.

\begin{definition}[Vacuum-pair invariant (D9(d))]
\label{def:vacuum-pair-invariant}
$G_{L} \times G_{R}$ acts on $\Vac \times \Vac$ componentwise. Carry
explicitly the \emph{transitivity hypothesis} (T): \emph{$G$ acts transitively
on $\Vac$}. The covariance requirement of \cref{def:covariant-vacuum} does
\emph{not} imply (T) --- the vacuum family may split into several $G$-orbits
--- so without (T) everything below holds per orbit. Under (T),
$\Vac \cong G/H_{\alpha}$ and
\begin{equation}
  \Vac \times \Vac \;\cong\; (G/H_{\alpha}) \times (G/H_{\alpha}),
\end{equation}
transitively, with the stabiliser of $(\alpha,\beta)$ equal to
$H_{\alpha} \times H_{\beta}$. This is \textbf{not} $(G \times G)/G_{\mathrm{diag}}$:
for $G = SU(2)$ with $H_{\alpha} = U(1)$ it is $S^{2} \times S^{2}$, of
dimension $4$, whereas $(G \times G)/G_{\mathrm{diag}} \cong SU(2)$ has
dimension $3$.

The complete invariant of a vacuum pair \emph{modulo the global (diagonal)
symmetry} is the \emph{double coset}
\begin{equation}
  \dcoset(\alpha_{L}, \alpha_{R})
  \; := \; H_{\alpha} \, g_{L}^{-1} g_{R} \, H_{\alpha}
  \;\in\; H_{\alpha} \backslash G / H_{\alpha},
  \label{eq:double-coset}
\end{equation}
where $\alpha_{L} = g_{L} \cdot \alpha$ and $\alpha_{R} = g_{R} \cdot \alpha$.
It is well defined --- changing $g_{L} \mapsto g_{L} h_{1}$ and
$g_{R} \mapsto g_{R} h_{2}$ conjugates within the double coset --- and
diagonal-invariant, since $g_{i} \mapsto h g_{i}$ cancels. For $G = SU(2)$
with $H_{\alpha} = U(1)$ it is the relative polar angle,
$\cos \theta = \hat{n}_{L} \cdot \hat{n}_{R} \in [-1,1]$.
\end{definition}

\begin{remark}[The corrected classification]
An earlier version asserted that the orbit of the coset space
$(G_{L} \times G_{R})/G_{\mathrm{diag}}$ is the set of vacuum pairs, and that
$[g_{L} g_{R}^{-1}]$ is the diagonal-invariant label. Both are false for
nonabelian $G$: the componentwise stabiliser is
$H_{\alpha} \times H_{\beta}$, not $G_{\mathrm{diag}}$, and
$g_{L} g_{R}^{-1} \mapsto h \, g_{L} g_{R}^{-1} h^{-1}$ under the diagonal
action. Both the double-coset invariant and the failure of
$g_{L} g_{R}^{-1}$ to be diagonal-invariant were verified numerically.
\end{remark}

\provenance{D9(d)}%
{\texttt{definitions.md} D9(d)}%
{\texttt{theory/checks/corner\_a\_check.py}}%
{\texttt{theory/verdicts/corner-a-r2.md} objection 5;
 \texttt{theory/verdicts/corner-a-r3.md}}

\subsection{Goldstone tensors and broken directions}

Differentiating the symmetry action on the vacuum tensor produces a tensor in
the excitation ansatz. This is the object whose zero-momentum limit is pure
gauge, and whose small-momentum behaviour is the Adler zero the campaign is
chasing.

\begin{definition}[Goldstone tensor, broken and unbroken directions (D11)]
\label{def:goldstone-tensor}
Assume the smoothness hypothesis (S) of \cref{def:smoothness}. Fix
$\alpha \in \Vac$, write $\hsub_{\alpha} := \mathrm{Lie}(H_{\alpha})$, and
choose an $\Ad(H_{\alpha})$-invariant complement $\msub_{\alpha}$ with
$\liealg = \hsub_{\alpha} \oplus \msub_{\alpha}$; the elements of
$\msub_{\alpha}$ are the \emph{broken directions}.

\emph{(a) Goldstone tensor.} For $\xi \in \liealg$,
\begin{equation}
  B_{\mathrm{G}}(\xi)
  \; := \; \frac{d}{d\varepsilon}
  \Bigl[\, \Utens\bigl(\exp(\varepsilon \xi)\bigr) A_{\alpha} \,\Bigr]
  \Big\vert_{\varepsilon = 0},
  \qquad\text{equivalently}\qquad
  B_{\mathrm{G}}(\xi)^{s} = \sum_{s'} q(\xi)_{s s'} \, A_{\alpha}^{s'} .
  \label{eq:goldstone-tensor}
\end{equation}

\emph{(b) Soft insertion.} In the sense of \cref{def:excitation-ansatz} and
\cref{def:modulated-charge},
\begin{equation}
  \ket{\Phi_{k}\bigl(B_{\mathrm{G}}(\xi)\bigr)} \;=\; Q_{k}(\xi) \triangleright \omega_{\alpha} :
\end{equation}
replacing $A_{\alpha}$ by $B_{\mathrm{G}}(\xi)$ at site $n$ is exactly the
action of the charge density $q_{n}(\xi)$ there.

\emph{(c) Goldstone count.} The number of independent Goldstone tensors at
$\alpha$ is
\begin{equation}
  \dim_{\C} \operatorname{span}_{\C}
  \bigl\{\, B_{\mathrm{G}}(\xi) \;:\; \xi \in \msub_{\alpha} \,\bigr\}
  \quad \text{modulo } \Nnull_{0}(\cdot),
\end{equation}
which may be strictly smaller than $\dim_{\R} \msub_{\alpha}$. The deficiency
is the \emph{type-B} phenomenon; the isotropic ferromagnet realises
$\dim_{\R} \msub_{\alpha} = 2$ broken directions but only one Goldstone mode.
\end{definition}

\provenance{D11}%
{\texttt{definitions.md} D11; the count and the ferromagnetic example at
 \texttt{theory/corner-a-goldstone.md} step (1.7)}%
{\texttt{theory/checks/corner\_a\_check.py}}%
{\texttt{theory/verdicts/corner-a-r3.md}}

\subsection{In which limits the ansatz gauge remainder vanishes}

The exact statement about null directions is the summation-by-parts identity,
whose remainder consists of exactly two boundary window vectors. Whether that
remainder can be dropped depends on the profile class, and different parts of
the campaign need different classes. The definition below states only what is
proved --- a bound on the remainder --- and then names the three regimes in
which that bound is used. Its type discipline is binding: quoting a fixed-$k$
identity with a decay hypothesis, or the reverse, is a defect.

\begin{definition}[The three regimes in which the gauge remainder is
negligible (D12)]
\label{def:gauge-limits}
Let $\Lambda_{L} := [-L, L]$ and
\begin{equation}
  \ket{\Phi_{f}^{\Lambda}(B)}
  \; := \; \sum_{n \in \Lambda} f(n) \,
  \ket{\psi_{\Lambda}(\ldots B@n \ldots)} .
\end{equation}
The summation-by-parts lemma shows that the gauge remainder
$\gaugerem_{\Lambda}[f, X]$ is a sum of exactly two boundary window vectors,
and gives the \textbf{upper bound}
\begin{equation}
  \bigl\Vert \gaugerem_{\Lambda}[f,X] \bigr\Vert
  \;\le\; 2 \, \Cbd \, \norm{X} \,
  \max\bigl(\abs{f(a)}, \abs{f(b)}\bigr),
  \label{eq:sbp-bound}
\end{equation}
with $\Cbd$ the window-norm constant \eqref{eq:cbd}. Everything below is a
consequence of this bound alone. Three regimes are used, and every statement
citing this definition must name which.

\emph{(a) Vanishing remainder for decaying profiles.} If $f \in c_{0}(\Z)$
then $\Vert \gaugerem_{\Lambda}[f,X] \Vert \to 0$ as
$\Lambda_{L} \nearrow \Z$. \textbf{This is a statement about the remainder
only.} It does \emph{not} assert that either side of the summation-by-parts
identity converges to a vector: for that one needs a summability class, as in
(a$'$), \emph{and} the split property, which remains a sketch. Measured decay
with the centred profile $(1+\abs{n-c})^{-3}$:
$\Vert \gaugerem_{\Lambda} \Vert / \Vert \text{bulk} \Vert
 = 3.4 \cdot 10^{-1},\, 5.7 \cdot 10^{-2},\, 8.5 \cdot 10^{-3},\,
 1.2 \cdot 10^{-3}$ for $L = 4, 8, 16, 32$.

\emph{(a$'$) Norm-convergent wave packets.} If in addition $f \in \ell^{1}(\Z)$
\textbf{and} $\Delta f \in \ell^{1}(\Z)$ --- that is, $f$ is summable of
bounded variation --- then both sides of the summation-by-parts identity are
absolutely convergent sums of window vectors of norm at most
$\Cbd \norm{X}$, uniformly in $\Lambda$; the identity then holds between the
two limits at fixed $\Lambda$ and remains remainderless as
$\Lambda \nearrow \Z$. A smooth compactly supported momentum packet
$f = \hat{\varphi}$ with $\varphi \in C_{c}^{\infty}((-\pi,\pi])$ has rapidly
decreasing $f$ and so lies in $\ell^{1} \cap BV$. The class $\Fell$ of
\cref{def:admissible-profiles} is exactly the bounded-variation condition.
\textbf{This class, not $c_{0}$, is the one to cite for wave-packet
statements.}

\emph{(b) $\delta$-normalised plane wave.} For $f(n) = e^{ikn}$ one has
$\abs{f} \equiv 1$, so \eqref{eq:sbp-bound} gives
$\Vert \gaugerem_{\Lambda} \Vert = O(1)$ \emph{uniformly in $\Lambda$} --- and
that is all that is needed:
\begin{equation}
  \bigl\Vert \, \abs{\Lambda}^{-1/2} \, \gaugerem_{\Lambda} \, \bigr\Vert
  \;=\; O\bigl(\abs{\Lambda}^{-1/2}\bigr) \;\longrightarrow\; 0 .
\end{equation}
No claim is made, or needed, about how the bulk term grows. Equivalently, per
site,
\begin{equation}
  \lim_{\Lambda \nearrow \Z} \abs{\Lambda}^{-1}
  \bra{\Phi_{k}^{\Lambda}(B')} O \, \ket{\gaugerem_{\Lambda}} \;=\; 0
\end{equation}
for every $O \in \Aloc$ and every fixed $B'$. Measured:
$\Vert \gaugerem_{\Lambda} \Vert \in [1.31,\, 2.03]$ for
$L = 4, \ldots, 32$.

\emph{Type discipline.} A fixed-$k$ identity may be quoted only with (b); a
$c_{0}$ or $\ell^{1}$ identity only with (a) or (a$'$), and then in the
real-space summation-by-parts form or as a Fourier superposition of the
fixed-$k$ identities --- never as a fixed-$k$ equation with a $c_{0}$
hypothesis. Statements quoted ``exactly'' without naming a regime are
defects.
\end{definition}

\begin{remark}[What this definition deliberately no longer says]
The r2 version over-quantified in three ways, all conceded. It asserted that
the identity ``holds exactly'' for every $f \in c_{0}$, although
$c_{0} \not\subset \ell^{1}$ and neither side need converge in norm --- take
$f(n) = (1+\abs{n})^{-1/4}$. It asserted a universal $\Theta(\abs{\Lambda}^{1/2})$
bulk growth, which is \emph{false}: for $\chi = 1$ one has
$A_{\alpha} X \propto A_{\alpha}$, so the plane-wave bulk sum is a bounded
geometric sum. And it attached a fixed-$k$ formula to a $c_{0}$ hypothesis,
even though a plane wave is not in $c_{0}$. What survives, and all that is
used, is the remainder bound \eqref{eq:sbp-bound}.
\end{remark}

\provenance{D12}%
{\texttt{definitions.md} D12 (rewritten in r3); remainder bound and
 summation-by-parts lemma at \texttt{theory/corner-a-goldstone.md} step (1.5);
 the sketch split property at \texttt{theory/corner-a.md} step (1.4),
 sub-step (2.9)}%
{\texttt{theory/checks/corner\_a\_check.py}, checks C4 (decay), C5 (plane-wave
 boundedness), C10 (the withdrawn bulk-growth assertion)}%
{\texttt{theory/verdicts/corner-a-r2.md} objection 4;
 \texttt{theory/verdicts/corner-a-r3.md}}
